% DeepPass68: general d-dimensional integral identities for loop calculations.

$d$ 维径向积分需要单位球面的面积。由定义
\begin{equation}
\Omega_d\equiv\int\dd\Omega_{d-1},
\label{eq:d-sphere-solid-angle-definition-dp68}
\end{equation}
可得
\begin{equation}
{
\Omega_d=\frac{2\pi^{d/2}}{\Gamma(d/2)}.}
\label{eq:d-sphere-solid-angle-dp21}
\end{equation}

把 Schwinger 参数式用于欧氏分母，先写成
\begin{equation}
\frac1{(\ell_E^2+\Delta)^\nu}
=\frac1{\Gamma(\nu)}
\int_0^\infty\dd s\,s^{\nu-1}\ee^{-s(\ell_E^2+\Delta)}.
\label{eq:euclidean-master-schwinger-dp68}
\end{equation}
在 $\Re\Delta>0$、$\Re\nu>d/2$ 的收敛域内先完成 $\ell_E$ 与 $s$ 积分，得到
\begin{proposition}[欧氏径向积分的 $\Gamma$ 表示]
先在质量平方为正、积分收敛的维数与指数区域成立；维数正则化将右侧亚纯延拓。有下列恒等式。
\begin{equation}
{
\int\frac{\dd^d \ell_E}{(2\pi)^d}
\frac1{(\ell_E^2+\Delta)^\nu}
=\frac1{(4\pi)^{d/2}}
\frac{\Gamma(\nu-d/2)}{\Gamma(\nu)}
\Delta^{d/2-\nu}.}
\label{eq:euclidean-master-integral-dp21}
\end{equation}
\end{proposition}
\begin{proof}
在正文给出的绝对收敛域内，可以交换 Schwinger 参数与圈动量积分。将\eqnrefs{eq:euclidean-master-schwinger-dp68} 代入，
\begin{align*}
 I_{d,\nu}(\Delta)
 &=\frac1{\Gamma(\nu)}
 \int_0^\infty\dd s\,s^{\nu-1}\ee^{-s\Delta}
 \int\frac{\dd^d \ell_E}{(2\pi)^d}\ee^{-s\ell_E^2}.
\end{align*}
用\eqnrefs{eq:d-dimensional-gaussian-r9} 完成圈动量积分，
\begin{align*}
 I_{d,\nu}(\Delta)
 &=\frac1{(4\pi)^{d/2}\Gamma(\nu)}
 \int_0^\infty\dd s\,
 s^{\nu-d/2-1}\ee^{-s\Delta}.
\end{align*}
再令 $u=s\Delta$，得到
\begin{equation*}
 \int_0^\infty\dd s\,s^{\nu-d/2-1}\ee^{-s\Delta}
 =\Delta^{d/2-\nu}\Gamma(\nu-d/2),
\end{equation*}
从而得到正文\eqnrefs{eq:euclidean-master-integral-dp21}。
\end{proof}
右端随后定义对 $d$、$\nu$ 与 $\Delta$ 的解析延拓。

四阶张量分子需要比\eqnrefs{eq:tensor-reduction-proof-r9} 多一个对称结构。完全对称且各向同性的秩四张量只能写成
\begin{equation}
\int\dd^d \ell\,\ell^\mu\ell^\nu\ell^\rho\ell^\sigma f(\ell^2)
=C_4\bigl(
\eta^{\mu\nu}\eta^{\rho\sigma}
+\eta^{\mu\rho}\eta^{\nu\sigma}
+\eta^{\mu\sigma}\eta^{\nu\rho}\bigr).
\label{eq:rank4-isotropic-ansatz-dp68}
\end{equation}
双重收缩固定系数后得到
\begin{equation}
{
\int\dd^d \ell\,
\ell^\mu\ell^\nu\ell^\rho\ell^\sigma f(\ell^2)
=\frac{\eta^{\mu\nu}\eta^{\rho\sigma}
+\eta^{\mu\rho}\eta^{\nu\sigma}
+\eta^{\mu\sigma}\eta^{\nu\rho}}
{d(d+2)}
\int\dd^d \ell\,(\ell^2)^2f(\ell^2).}
\label{eq:rank4-tensor-reduction-dp21}
\end{equation}

在 $d=4-2\epsilon$ 附近，圈积分中经常出现同一个无量纲组合
\begin{equation}
\mathcal U_\epsilon(\Delta)
\equiv\Gamma(\epsilon)(4\pi)^\epsilon
\left(\frac{\Delta}{\mu^2}\right)^{-\epsilon},
\label{eq:dimreg-universal-factor-dp68}
\end{equation}
其中 $\mu$ 是保持重整化参数具有固定工程量纲所引入的尺度。小 $\epsilon$ 展开为
\begin{proposition}[维数正则化因子的 Laurent 展开]
在正文的重整化尺度与 $\overline{\mathrm{MS}}$ 约定下，取维数偏移趋于零并固定其余量。有下列恒等式。
\begin{equation}
{
\mathcal U_\epsilon(\Delta)
=\frac1\epsilon-\gamma_E+\ln4\pi
-\ln\frac{\Delta}{\mu^2}+\order(\epsilon).}
\label{eq:dimreg-universal-expansion-dp68}
\end{equation}
\end{proposition}
\begin{proof}
维数正规化中的方案常数来自三个普通小参数展开。首先由 $\Gamma(1+\epsilon)=\epsilon\Gamma(\epsilon)$ 和
\begin{equation*}
 \Gamma(1+\epsilon)=1-\gamma_E\epsilon+\order(\epsilon^2)
\end{equation*}
得到
\begin{equation*}
 \Gamma(\epsilon)=\frac1\epsilon-\gamma_E+\order(\epsilon).
\end{equation*}
另外
\begin{align*}
 (4\pi)^\epsilon
 &=1+\epsilon\ln4\pi+\order(\epsilon^2),\\
 \left(\frac{\Delta}{\mu^2}\right)^{-\epsilon}
 &=1-\epsilon\ln\frac{\Delta}{\mu^2}+\order(\epsilon^2).
\end{align*}
三者相乘并保留有限阶，得到
\begin{equation*}
 \mathcal U_\epsilon(\Delta)
 =\frac1\epsilon-\gamma_E+\ln4\pi
 -\ln\frac{\Delta}{\mu^2}+\order(\epsilon),
\end{equation*}
即正文\eqnrefs{eq:dimreg-universal-expansion-dp68}。$\overline{\rm MS}$ 方案定义为在这个通用展开中一并减去 $1/\epsilon-\gamma_E+\ln4\pi$；这一步定义了 $\overline{\rm MS}$ 的减除方案；母积分本身仍由前面的维数正则化恒等式给出。
\end{proof}
因此 $\overline{\rm MS}$ 中反复出现的 $1/\epsilon-\gamma_E+\ln4\pi$ 只是这个通用展开的方案常数部分。
